\section{Introduction}
\label{sec:intro}

Training fabrics for large machine-learning models spray packets. Instead of pinning a flow to one path, the leaf switch sends consecutive packets of the same flow to different spines, which balances load across a Clos fabric and is now the default of the Ultra Ethernet transport~\cite{spraying,uec,reps}. Spraying also changes what a failure looks like. A link that drops one packet in a thousand, or one in ten thousand, does not take a flow down. It slows every flow a little, because every flow crosses it. Such gray failures are common in production fabrics and hard to attribute~\cite{grayfailure,corropt}, and they can be costly for synchronous training jobs whose collective operations wait for the slowest packet~\cite{aegis}. We call a link that silently drops a small fraction of its packets a grayhole, and we call the problem of naming the directed link that carries it grayhole localization.

Grayhole localization on a sprayed fabric is challenging, due to two reasons. First, the evidence is divided. A flow that used to concentrate its loss on one path now spreads it across $k$ spines. Consequently, the per-path deficit a passive observer can see is $1/k$ of the fault, and it shrinks further as the loss rate falls, while the sampling noise of spraying, i.e., the variation in how many packets each spine happens to receive, does not shrink at all. Second, the evidence is aliased. A packet from one leaf to another crosses two directed links, an uplink into a spine and a downlink out of it. An arrival counter at the destination sees only the product of the two survival rates, and thus, a single vantage that counts arrivals can name a spine but cannot, by itself, tell the uplink from the downlink.

The two most recent detectors built for sprayed fabrics accept both limits. SprayCheck~\cite{spraycheck} applies a Z-test to per-spine arrival counts and resolves a directed link by intersecting reports from several leaves. FlowPulse~\cite{flowpulse} compares each ingress port's load against a learned baseline and applies a per-sender rule. Both are passive: they add nothing to the wire, they run on counters the switch already has, and they can be deployed today. However, a test on arrival counts needs enough packets for a deficit of fraction $p$ to stand above spray noise, and that number grows at least as $1/p$. Both papers also state that one report can only flag the uplink and downlink pair~\cite{spraycheck,flowpulse}.

An in-fabric witness can remove both limits by construction. The sending leaf stamps every packet with a per-directed-link sequence number. The receiving switch keeps, per directed link, the highest sequence number seen and the number of packets seen. The difference of the two deltas is then the loss on that directed link and nothing else. This primitive is not new. NetSeer~\cite{netseer}, LinkGuardian~\cite{linkguardian}, LossRadar~\cite{lossradar}, and the link-layer retry of Ultra Ethernet~\cite{uec} all place a per-link sequence witness in the data plane. What none of them reports is what the witness buys, and costs, on a sprayed fabric at the loss rates where the passive state of the art fails. That measurement, head-to-head against the passive detectors and on real silicon, is the contribution of this paper. We implement the witness as a receiver ledger on an Intel Tofino switch~\cite{tofino}, with a two-byte stamp and two registers per directed link, and we treat it as the measured arm of a comparison rather than as a new design.

To make the comparison fair, we replay SprayCheck and FlowPulse from their own papers, with their own thresholds and their own localization rules. All three detectors read one shared stream of sprayed traffic, and each sees it only through the counters its own switch would have. We run fifty seeds per loss rate from 1.5 percent down to $10^{-4}$, and we report false positives next to every action rate. Where a passive detector ties or beats the witness, we report it as plainly as where it does not. Our contributions are as follows.

\begin{itemize}
\item \textbf{Measures the detection cost-scaling separation.} The witness detects a single-link grayhole on the first epoch after onset at every loss rate. Its median cost is a flat 22.0 million fleet packets from 1.5 percent to $10^{-4}$, with an action rate of 1.00 and no false positive. SprayCheck's median cost grows from 24 million to 114 million packets as the loss rate falls to $10^{-3}$, its action rate falls from 1.00 to 0.78, and at $10^{-4}$ it detects nothing within the same 160-million-packet budget. FlowPulse fails one order of magnitude earlier. We show that this is a scaling law of any arrival-count test and not an artifact of the budget (Section~\ref{sec:detection}).
\item \textbf{Measures directed-link localization under spraying.} The witness names the exact directed link, and only it, at every loss rate and for both uplink and downlink faults. SprayCheck ties it at 1.0 to 1.5 percent loss, which we report as a null result, and degrades to the ambiguous uplink and downlink pair at $10^{-3}$. FlowPulse cannot detect an uplink fault at any rate once several senders share the port (Section~\ref{sec:localization}).
\item \textbf{Measures the ledger on silicon.} On a Tofino switch the ledger recovers injected loss exactly from $10^{-2}$ down to $10^{-4}$, i.e., six drops in sixty thousand packets, with zero false positives on more than forty thousand clean packets. The result holds for scheduled drops and for random per-packet drops, and a downlink grayhole is attributed to the downlink alone with zero false attribution on seven active sublinks (Section~\ref{sec:silicon}).
\item \textbf{Bounds the claim with a correlated-fault gate.} Under two or three simultaneous independent faults the witness's advantage widens, because the passive intersection produces spurious links. Under a fleet-wide shift in loss with no single culprit the witness fails and flags all 64 directed links, and SprayCheck's relative test is the better tool. We report both, and we scope every other claim to independent per-directed-link faults (Section~\ref{sec:robustness}).
\item \textbf{Quantifies the cost.} The two-byte witness adds about 0.14 percent to a 1400-byte payload and six bytes of register state per directed link. It also costs one ingress pipeline stage, which brings the program to the twelve-stage ceiling of the Tofino~1 (Section~\ref{sec:cost}).
\end{itemize}

To the best of our knowledge, this is the first head-to-head measurement of an in-fabric per-directed-link witness against the passive sprayed-fabric detectors, on the loss regime all three target, and on silicon. We claim no novelty for the witness itself. The per-link sequence stamp is the primitive of NetSeer and LinkGuardian, and the localization advantage we measure is a property of the information the witness holds rather than of a better inference. What the paper adds is the measurement, its cost, and its boundary.

The rest of the paper is organized as follows. Section~\ref{sec:background} describes sprayed fabrics, grayholes, the fault model, and the objectives we evaluate. Section~\ref{sec:ledger} describes the receiver ledger as implemented. Section~\ref{sec:method} describes the replay harness, the two passive baselines, and the silicon testbed. Sections~\ref{sec:detection} to~\ref{sec:cost} report detection cost, localization, the silicon measurements, the correlated-fault gate, and the cost. Section~\ref{sec:discussion} discusses limitations, Section~\ref{sec:related} positions the work against prior systems, and Section~\ref{sec:conclusion} concludes.
